\begin{proposition*}[Degree one and scalar operators]
On $V\ne\{0\}$, $\deg m_T=1$ iff $T=cI$ for some $c\in F$; then $m_T=x-c$.
\end{proposition*}
\begin{prob}\begin{proof}\begin{steps}
\ii A monic degree-one relation $m_T=x-c$ says $T-cI=0$.
\ii Conversely, $x-c$ is zero on $cI$, while a nonzero constant gives a nonzero scalar multiple of $I$ on a nonzero space. Thus no degree-zero polynomial works.
\end{steps}\end{proof}\end{prob}
\begin{example*}[A vector relation that is not global]
Over $\mathbb R$, let $A=\begin{pmatrix}2&0&0\\1&2&0\\0&0&3\end{pmatrix}$.
\begin{steps}
\ii $(A-2I)e_1=e_2$ and $(A-2I)e_2=0$. Thus $p_{e_1}=(x-2)^2$: the quadratic works, and $e_1,Ae_1=(2,1,0)$ are independent, excluding a linear relation. Also $p_{e_2}=x-2$ and $p_{e_3}=x-3$ by their eigenvector equations.
\ii The basis criterion gives $m_A=(x-2)^2(x-3)$. Using only $e_1$ would have missed the required root $3$. To compute the least common multiple here, keep each distinct factor with its largest required exponent; the factor identity and successive cancellation justify this for split polynomials.
\end{steps}\end{example*}
\begin{proposition*}[Direct summands and matrix blocks]
If $V=U_1\oplus\cdots\oplus U_r$ with each $U_j$ invariant, then $m_T=\operatorname{lcm}(m_{T|_{U_1}},\ldots,m_{T|_{U_r}})$. Thus a block diagonal matrix, with square blocks on the diagonal and zero off-diagonal blocks, has the least common multiple of the block minimal polynomials.
For $A=\begin{pmatrix}B&E\\0&C\end{pmatrix}$, with $B,C$ square, put $r=m_B$, $s=m_C$. Then
\[\operatorname{lcm}(r,s)\mid m_A\mid rs.\]
If $\operatorname{lcm}(r,s)=rs$, both bounds give $m_A=rs$.
\end{proposition*}
\begin{prob}\begin{proof}\begin{steps}
\ii For $v=\sum_j u_j$, $q(T)v=\sum_j q(T|_{U_j})u_j$. Thus $q(T)=0$ iff $q(T|_{U_j})=0$ on every summand, iff every block minimal polynomial divides $q$. The least-common-multiple definition proves the first assertion.
\ii In the upper-block case, the first coordinates form an invariant subspace on which $A$ acts by $B$; the quotient, identified by the remaining coordinates, acts by $C$. Relations on invariant parts give $r,s\mid m_A$ and $m_A\mid rs$. The defining property of the least common multiple gives the lower bound.
\end{steps}\end{proof}\end{prob}
\begin{example*}[Midterm 2 Review, Problem 5: eigenvalue and relation parts]
Let $B=\begin{pmatrix}3&1&1\\0&3&-1\\0&0&1\end{pmatrix}$ over $\mathbb R$.
\begin{steps}
\ii The top-left block has minimal polynomial $(x-3)^2$ because subtracting $3I$ gives a nonzero matrix with square zero. The last block has minimal polynomial $x-1$. These have no shared root, so their least common multiple is their product; the block bounds give $m_B=(x-3)^2(x-1)$.
\ii Triangularity gives eigenvalues $3,1$. For $\lambda=3$, the equations are $y+z=0,-z=0,-2z=0$, giving $E(3,B)=\mathbb Re_1$. For $\lambda=1$, they are $2x+y+z=0,2y-z=0$, giving $E(1,B)=\mathbb R(-3,2,4)$.
\ii The eigenspace dimensions total $2<3$, so $B$ is not diagonalizable.
\end{steps}\end{example*}
\begin{proposition*}[Two-by-two relation; LADR 5B Exercise 11]
For $A=\begin{pmatrix}a&b\\c&d\end{pmatrix}$,
\[A^2-(a+d)A+(ad-bc)I=0.\]
If $A$ is nonscalar, this monic quadratic is $m_A$; if $A=\gamma I$, then $m_A=x-\gamma$.
\end{proposition*}
\begin{prob}\begin{proof}\begin{steps}
\ii Multiplication gives $A^2=\left(\begin{smallmatrix}a^2+bc&b(a+d)\\c(a+d)&d^2+bc\end{smallmatrix}\right)$. Subtracting $(a+d)A$ leaves $(bc-ad)I$, proving the relation directly in entries.
\ii A nonscalar matrix has no degree-one relation by the scalar criterion, so this degree-two relation is minimal. This uses no general determinant theorem.
\end{steps}\end{proof}\end{prob}
\begin{example*}[A relation from a rank-one product]
For nonzero columns $u,v\in F^n$, $A=uv^{\mathsf T}$ has every column a multiple of $u$, and
$A^2=u(v^{\mathsf T}u)v^{\mathsf T}=(v^{\mathsf T}u)A$.
For $u=(1,2)^{\mathsf T}$, $v=(3,0)^{\mathsf T}$ over $\mathbb R$, this gives
$A=\left(\begin{smallmatrix}3&0\\6&0\end{smallmatrix}\right)$ and $A^2=3A$.
Since $A$ is nonscalar, $m_A=x(x-3)$. The factorization of the matrix itself suggested the relation.
\end{example*}
\begin{proposition*}[Reduction to low-degree polynomials]
Every polynomial value $f(T)$ has a unique expression $r(T)$ with $r=0$ or $\deg r<\deg m_T$. In particular, divide $x^k=m_Th+r$ to compute $T^k=r(T)$.
\end{proposition*}
\begin{prob}\begin{proof}\begin{steps}
\ii Polynomial division gives $f=m_Th+r$, so $f(T)=r(T)$ by $m_T(T)=0$.
\ii If two such remainders give the same operator, their difference is zero on $T$ and has degree below $\deg m_T$ if nonzero. Minimality excludes that possibility, proving uniqueness. On the zero space the only allowed remainder is zero.
\end{steps}\end{proof}\end{prob}
\method{Compute a high power from a remainder}
\begin{steps}
\ii For $A=\left(\begin{smallmatrix}2&1\\1&2\end{smallmatrix}\right)$ over $\mathbb R$, $m_A=(x-1)(x-3)$. The remainder of $x^k$ has form $a x+b$. Evaluating the division identity at $1,3$ gives $a+b=1$ and $3a+b=3^k$.
\ii Hence $a=(3^k-1)/2$, $b=(3-3^k)/2$. For $k=4$, $r=40x-39$ and $A^4=40A-39I=\left(\begin{smallmatrix}41&40\\40&41\end{smallmatrix}\right)$.
\end{steps}
\begin{proposition*}[Invertibility and a polynomial inverse; LADR 5.32]
$T$ is invertible iff $m_T(0)\ne0$. For $\dim V>0$, write $m_T=\sum_{j=0}^d a_jx^j$, $a_d=1$. When $a_0\ne0$,
\[T^{-1}=-a_0^{-1}\sum_{j=1}^d a_jT^{j-1}.\]
On the zero space $m_T=1$ and its unique operator is its own inverse.
\end{proposition*}
\begin{prob}\begin{proof}\begin{steps}
\ii By Proposition 66, $m_T(0)=0$ iff $0$ is an eigenvalue, iff $T$ is not invertible by the shifted-kernel criterion (Week 6).
\ii Rearrange $\sum_{j=1}^d a_jT^j=-a_0I$, and factor out $T$. Dividing by $-a_0$ gives a right inverse; it is a polynomial in $T$, so it commutes with $T$ and is also a left inverse.
\end{steps}\end{proof}\end{prob}
\method{Compute an inverse from a relation}
\begin{steps}
\ii For $A=\left(\begin{smallmatrix}2&1\\0&3\end{smallmatrix}\right)$ over $\mathbb R$, multiplication gives $A^2-5A+6I=0$. Thus $A(5I-A)=6I$.
\ii Therefore $A^{-1}=(5I-A)/6=\left(\begin{smallmatrix}1/2&-1/6\\0&1/3\end{smallmatrix}\right)$. Multiplying by $A$ gives $\left(\begin{smallmatrix}1&0\\0&1\end{smallmatrix}\right)$: the off-diagonal entry is $-2/6+1/3=0$.
\end{steps}
\begin{prob}\textbf{Changing coordinates or transposing preserves relations.}
\begin{steps}
\ii For invertible $S$, induction gives $(S^{-1}AS)^k=S^{-1}A^kS$, so $p(S^{-1}AS)=S^{-1}p(A)S$. This is zero iff $p(A)=0$, hence $m_{S^{-1}AS}=m_A$.
\ii $(A^{\mathsf T})^k=(A^k)^{\mathsf T}$, so $p(A^{\mathsf T})=p(A)^{\mathsf T}$ is zero iff $p(A)=0$, giving $m_{A^{\mathsf T}}=m_A$.
\end{steps}\end{prob}
\begin{prob}\textbf{Divisors of split polynomials.}
A monic divisor of $\prod_j(x-\lambda_j)$ is a product of some of these factors, with no factor used more often than it appears in the product.
\begin{steps}
\ii Write $ph=\prod_j(x-\lambda_j)$. At $\lambda_1$, either $p(\lambda_1)=0$ or $h(\lambda_1)=0$. The root factor identity removes $x-\lambda_1$ from that factor; cancel it from both sides.
\ii Continue with the remaining linear factors. At the end $p$ is reduced to a constant; its removed factors give its factorization, and monicity makes the constant $1$. Distinct factors in the original product therefore remain distinct in a divisor.
\end{steps}\end{prob}
\begin{proposition*}[Compare triangular and diagonal forms; LADR 5.40, 5.44]
$T$ has an upper-triangular matrix over $F$ iff $m_T$ splits into linear factors over $F$; the factors may repeat. The empty basis and polynomial $1$ handle the zero space.
\end{proposition*}
\begin{prob}\begin{proof}\begin{steps}
\ii In a triangular basis with diagonal $\lambda_j$, put $U_j=\operatorname{span}(v_1,\ldots,v_j)$, $U_0=\{0\}$. The column condition gives $(T-\lambda_jI)U_j\subseteq U_{j-1}$. Induction on $j$ shows $\prod_{i=1}^j(T-\lambda_iI)$ is zero on $U_j$: the last factor sends into $U_{j-1}$ and the preceding product is zero there. Thus $\prod_{i=1}^n(x-\lambda_i)$ is a global relation. Its divisor $m_T$ splits by the preceding argument.
\ii Conversely induct on $\dim V$, starting with the empty basis. For nonzero $V$ with split $m_T$, a root $\lambda$ is an eigenvalue by Proposition 66. The proper invariant subspace $U=\im(T-\lambda I)$ has $m_{T|_U}\mid m_T$, hence its polynomial splits too. Induction supplies a triangular basis of $U$.
\ii Extend it by $w_j$ to a basis of $V$. Each new vector satisfies $Tw_j=(T-\lambda I)w_j+\lambda w_j\in U+Fw_j$. Along with the initial-span condition inside $U$, this is precisely Proposition 59's triangular column condition.
\end{steps}\end{proof}\end{prob}
